%!TEX root = ../thesis.tex
%*******************************************************************************
%******************************** Eighth Chapter *******************************
%*******************************************************************************

\chapter{Conclusion and Future Work}

\section{Summary}
This thesis has defended the principle of structure-aligned learning (Section~\ref{sec:unifying_argument}): reliable constraint satisfaction requires aligning learning with explicit task-relevant structure. Depending on the task, this alignment may modify the model architecture or representation, introduce learning objectives or explicit supervision, alter the decoding procedure, or enable data construction and training to scale. The principle was examined in visual understanding and visual generation, two settings with different outputs and mechanisms but a shared need to prevent statistically convenient shortcuts from displacing task-defining relationships.

For visual understanding, the task-relevant structure is a set of pixel-grounded entity--relation triplets. This thesis focused on panoptic scene graph generation (PSG)~\citep{yang2022panoptic}, which represents both ``things'' and ``stuff'' as graph entities. Chapters~3--5 developed three complementary advances. \textbf{HiLo}~\citep{zhou2023hilo} improves long-tail relation recognition under semantic overlap through frequency-specialized branches and cross-branch consistency. \textbf{VLPrompt}~\citep{zhou2025vlprompt} aligns context-rich language cues from large language models with grounded visual pair features. \textbf{OpenPSG}~\citep{zhou2025openpsg} moves PSG toward open-world prediction through open-set panoptic segmentation, efficient pair filtering, and autoregressive multimodal relation decoding.

For visual generation, the task-relevant structure comprises pose, spatial correspondence, and identity consistency over time. This thesis studied reference-conditioned image and video synthesis with diffusion models~\citep{ho2020denoising, rombach2022high}, where these relationships must be preserved without compromising realism. \textbf{Leffa}~\citep{zhou2025learning} regularizes attention as a reference--target correspondence field, guiding the denoiser toward the correct source regions and reducing fine-grained detail distortion. \textbf{Saber}~\citep{zhou2025scaling} uses masked training, tailored attention, and mask augmentation to learn reference-aware temporal identity from ordinary video--text pairs, without expensive reference image--video--text triplets.

The five contributions instantiate the same principle without sharing a single architecture. Chapters~3--5 make relational structure operational through branch specialization and consistency, grounded vision--language fusion, and pair-aware generative decoding. Chapters~6--7 act on spatial and temporal correspondence through attention regularization and masked reference simulation. These choices address different shortcuts: dominance by frequent predicates, reliance on ungrounded linguistic priors, restriction to a fixed vocabulary, diffuse reference attention, and adjacent-frame copying. The common result is that task-defining relationships exert a direct influence on learning instead of remaining implicit in the input or label space.

\noindent \textbf{Contributions at a Glance.}
Across Chapters~3--7, the contributions of this thesis can be summarized as five algorithmic steps toward more dependable systems:
\begin{itemize}
    \item \textbf{HiLo: multi-granular long-tail PSG.} HiLo addresses a mismatch between annotation and semantics, since relation labels are both imbalanced and overlapping. It specializes modeling capacity for head and tail predicates and enforces consistency across granularities, which improves tail relation prediction without sacrificing overall structure quality~\citep{zhou2023hilo}.
    \item \textbf{VLPrompt: language priors for relation prediction.} VLPrompt uses large language models to provide context-rich relation descriptions that complement visual evidence. It does not replace vision with language. Instead, it treats language as an auxiliary signal and learns when to trust it, which improves generalization for rare predicates and ambiguous contexts~\citep{zhou2025vlprompt}.
    \item \textbf{OpenPSG: open-world grounded relation prediction.} OpenPSG extends PSG beyond a fixed predicate vocabulary by generating grounded relation phrases with multimodal models and by using efficiency-aware pair reasoning. This changes the interface from closed-set classification to open-world grounded prediction~\citep{zhou2025openpsg}.
    \item \textbf{Leffa: attention regularization for reference preservation.} Leffa interprets cross-reference attention as a soft correspondence field and regularizes it during diffusion training through learnable flow fields. This improves the preservation of textures, logos, and fine patterns under large pose or appearance changes, without adding inference-time modules~\citep{zhou2025learning}.
    \item \textbf{Saber: scalable reference-to-video via simulated references.} Saber removes the reliance on expensive reference image--video--text triplets by training on ordinary video--text pairs and simulating reference conditioning through masked training, attention design, and augmentation. The resulting model preserves identity over time and generalizes across different numbers of reference images~\citep{zhou2025scaling}.
\end{itemize}
\noindent The first three contributions improve the learning and decoding of pixel-grounded relational structure; the last two improve the learning of spatial and temporal reference structure. Together, they instantiate structure-aligned learning through model architecture and representation, learning objectives and explicit supervision, decoding procedures, and scalable data construction and training strategies.

\section{Lessons Learned}
Four consequences of the central claim recur across the chapters. They concern grounding, long-tail learning, external knowledge, and reference-conditioned generation, but each follows from the same principle: learning must be aligned with the structure that defines success.

\noindent \textbf{(i) Task-Relevant Structure Must Be Explicit and Grounded.}
Object categories alone do not express how a scene is organized. A subject--object--relation triplet exposes compositional structure, while panoptic masks ground its entities at the pixel level~\citep{cheng2022visual, zhou2023hilo}. Both components are necessary: relations without grounding can be plausible but visually unsupported, whereas masks without relations omit the interactions that distinguish a structured scene description from a list of regions. Semantic granularity is also part of this structure. Several predicates may be valid for the same pair at different levels of specificity, so collapsing them as label noise removes information that the task is intended to recover.

\noindent \textbf{(ii) Learning Must Reflect Structure, Not Only Marginal Frequency.}
Long-tail performance depends on how modeling capacity and gradients are allocated, not only on sample counts. Relation prediction is especially vulnerable to shortcuts because object labels and co-occurrence statistics predict many head predicates without sufficient visual evidence. HiLo makes the head--tail organization operational through specialized branches and cross-branch consistency, preventing frequent gradients from erasing rarer but semantically informative relations~\citep{tang2020unbiased, zellers2018neural, zhou2023hilo}. Evaluation must reflect the same structure: recall and mean recall expose complementary aspects of performance, and neither is adequate in isolation~\citep{zellers2018neural}.

\noindent \textbf{(iii) Language-Derived Structure Must Be Aligned with Visual Evidence.}
Language models and vision--language pretraining provide relational knowledge beyond the statistics available in a task dataset. Used naively, however, this knowledge introduces hallucinations or overly generic predictions~\citep{radford2021learning, min2023recent, zhang2024vision}. An ungrounded prior often shifts the output toward the typical case, reinforcing rather than correcting frequency bias. Language is therefore most effective as a complement to pair-specific visual evidence, with learned fusion determining when it is informative~\citep{zhou2025vlprompt}. In open-world settings, linguistic flexibility must likewise be balanced against grounding, calibration, and abstention~\citep{wu2024towards, zhang2024vision}.

\noindent \textbf{(iv) Reference Structure Must Shape the Learning Objective.}
Additional conditions such as text, masks, depth, or references do not by themselves ensure faithful generation~\citep{zhang2023adding, xie2023boxdiff}. The model must bind a condition to the relevant target region and preserve the required relationship along the denoising trajectory and over time. Leffa makes spatial correspondence operational through attention-flow regularization. Saber makes temporal identity operational through masked training that requires reference-aware reconstruction while disrupting a direct copy shortcut~\citep{zhou2025learning, zhou2025scaling}. These results motivate objectives that target preservation and correspondence directly, while scaling beyond highly curated datasets.

\section{Limitations}
\label{sec:thesis_limitations}
Each technical chapter states the limitations of its own method (Sections~\ref{hilo:sec:limitations}, \ref{sec:limitation}, \ref{openpsg:sec:limitations}, \ref{leffa:sec:limitations}, and~\ref{saber:sec:limitations}). This section has a different purpose. It identifies the limitations that recur across chapters, because those follow from the argument of this thesis rather than from any particular design, and are therefore harder to remove.

\noindent \textbf{(i) This thesis measures agreement with an annotation, not correctness.}
This is the most pervasive limitation of the work. It takes a different form in each part but remains the same problem. In Part~I, PSG annotations are incomplete. R@$K$ and mR@$K$ score predictions against a label set that is known to omit valid relations, so a correct but unannotated relation is penalized twice: it receives no credit, and it also displaces an annotated relation from the top-$K$ list (Sec.~\ref{hilo:sec:limitations}). For a thesis about the long tail this is not a marginal effect, because the omissions concentrate in exactly the rare and specific predicates that the work targets. In Part~II, the mismatch is between the property we claim and the property we measure. FID is a distributional statistic that is largely blind to whether a logo is legible, and OpenS2V-Eval aggregates eight model-based scores whose weighting we did not choose (Sections~\ref{leffa:sec:limitations} and~\ref{saber:sec:limitations}).

The consequence is that the improvements reported in this thesis are improvements in agreement with a benchmark, and the gap between that and correctness is unmeasured. The human studies in Chapter~6 and the failure analysis in Chapter~4 narrow this gap but do not close it. Stated in its strongest form: no experiment in this thesis can distinguish a model that has become better at the task from a model that has become better at matching the conventions of the annotators.

\noindent \textbf{(ii) The structure made operational inherits the biases of its definition.}
The choice of task-relevant structure determines which distinctions a model can learn, and every structure used here reflects prior design decisions. The $56$ PSG predicates were fixed by the dataset authors over COCO images. The open-vocabulary reach of Chapter~5 is bounded in practice by a frozen BLIP-2 decoder. The pose representation of Chapter~6 assumes a skeleton suited to upright and largely unoccluded human bodies. Relations that are visually subtle, culturally specific, or absent from these representations cannot be recovered by improving the learning procedure alone. Explicit structure makes these biases inspectable; it does not make the system neutral.

\noindent \textbf{(iii) Confident wrong answers, and the absence of abstention.}
None of the five systems can say that it does not know. The concern is sharpest where language models are involved. In Chapter~4, LLM descriptions are unverified against the image and biased toward stereotypical interactions, so the hallucination risk aligns with the head of the distribution, which is the bias the chapter aims to reduce. In Chapter~5, a frozen decoder produces a fluent relation phrase for any pair that survives filtering, so a false positive from the pair filter becomes a confident and wrong edge in the graph. In Part~II, a generator asked for a detail that the reference does not contain will synthesize one rather than decline. Learned gating in Chapter~4 and pair filtering in Chapter~5 are partial defenses, but both are trained on the same long-tailed data whose tail behavior they are meant to calibrate. This thesis does not address calibration or abstention, and for open-world deployment neither is optional.

\noindent \textbf{(iv) Errors cascade, and the weakest component is usually not the contributed one.}
Every system here is a pipeline built on frozen or pretrained components, and in each case the dominant failure mode originates upstream of the contribution. Two adjacent people merged into a single panoptic mask make every relation involving them wrong, however good the relation decoder is (Chapters~3--5). An inaccurate garment segmentation changes which pixels are synthesized, however well correspondence is resolved (Chapter~6). The same dependency reappears at inference in the foreground segmenter of Chapter~7. This has a methodological implication that we do not fully honor. Ablations that hold the upstream components fixed measure the contribution under favorable conditions, so the reported gains are upper bounds on what a deployed system with realistic upstream error rates would show.

\noindent \textbf{(v) Scope of the evidence.}
The empirical basis is narrower than the argument. Part~I is evaluated on PSG and VG-150, two benchmarks that share much of the same photographic and annotation practice. Part~II is evaluated on three person-image datasets with broadly consistent capture conditions, and on a single video benchmark. Nothing here tests robustness to domain shift, such as robotics, medical imagery, or non-Western visual culture. The long-tail claims in particular are claims about the tail of these annotations rather than about the tail of the visual world.

\noindent \textbf{(vi) Cost has not been reduced, only relocated.}
This thesis repeatedly trades one cost for another. HiLo doubles decoder computation in order to separate head and tail behavior. VLPrompt removes per-image LLM cost by pre-extracting class-level descriptions, at the price of a language signal that cannot respond to the specific instance. Pair filtering in OpenPSG reduces the quadratic cost but does not remove it, and every surviving pair still requires a large-decoder forward pass. Leffa adds no inference cost but requires a multi-stage training schedule with no principled criterion for when to switch the loss on. Saber removes the need for curated triplet data but still finetunes a 14B model and runs $50$ denoising steps at inference. Concretely, in Chapter~6 we train the try-on model at $512 \times 384$ for $10\text{k}$ steps, at $1024 \times 768$ for $36\text{k}$ steps, and then finetune for a further $10\text{k}$ steps with the \leffa{} loss~\citep{rombach2022high, zhou2025learning}. In Chapter~7, Saber is finetuned from Wan2.1-14B and depends on a foreground segmenter at test time~\citep{wan2025wan, zhou2025scaling}. Structure-aligned learning is therefore not free, and none of the mechanisms proposed here is cheap enough to adopt without deliberation.

\noindent \textbf{(vii) The principle is instantiated separately rather than integrated end to end.}
The final limitation concerns the thesis as a whole. Section~\ref{sec:unifying_argument} identifies a common design principle across understanding and generation, but no chapter connects them in one system. A predicted scene graph is not used to condition a generator or to evaluate its output. Such a system would test whether structure alignment can remain effective across the boundary between analysis and synthesis. Section~\ref{sec:future_structure_conditioned} outlines this extension.

\section{Future Directions}
Each of the directions below is described at the level of a research program rather than an open problem. Each states concrete research questions, the methodological opportunities that make them tractable now, and the evaluation challenges that must be solved before progress can be measured. They are ordered by how directly they follow from the argument of this thesis. The first would test that argument in its strongest form, and the remaining ones address the limitations listed in Section~\ref{sec:thesis_limitations}.

\subsection{Structure-Conditioned and Structure-Evaluated Generation}
\label{sec:future_structure_conditioned}
This thesis applies structure-aligned learning separately to visual understanding and visual generation. A stronger test would connect the two: the scene structure recovered by an understanding system would directly condition a generator and evaluate its output. No chapter here uses a scene graph in either role. Closing this gap is therefore the most direct extension of the thesis and constitutes a research program rather than a single experiment.

\noindent \textbf{Research questions.}
First, is a scene graph a better generation interface than a caption, and for which properties? A graph specifies which entities exist, how they are related, and, because it is pixel-grounded, where they are. These are exactly the properties that captions specify poorly. A graph is also a lossy description of an image, since it says nothing about style, lighting, or the appearance of an entity beyond its category. The question is therefore not whether graphs replace prompts, but which control properties each interface suits, and whether a hybrid interface inherits the strengths of both or the ambiguities of both.

Second, how should a graph be injected into a diffusion model without collapsing into a bag of nouns? Text encoders are known to lose relational structure, which is why a prompt for a cat to the left of a dog often produces the reverse arrangement. Injecting a graph as serialized text would reproduce this failure. The question is whether relational structure survives conditioning better when the graph is encoded as a typed set of entity tokens with explicit edge tokens, when edges constrain attention directly between the token groups of the related entities, or when the mask geometry of each panoptic entity acts as a spatial prior for where its tokens may attend.

Third, what should a generator do with an inconsistent or underspecified graph? Real graphs produced by a PSG model contain contradictions and omissions. A system that fails badly on such input is unusable, so handling partial specifications is part of the problem rather than an implementation detail.

Fourth, does a graph-conditioned generator inherit the long-tail bias of the graphs it was trained on? If the training graphs come from a model with the frequency bias documented in Chapter~3, the generator will render \textit{on} more reliably than \textit{riding}. This is a concrete and testable prediction, and confirming it would show that structural bias propagates across the boundary between understanding and generation.

\noindent \textbf{Methodological opportunities.}
The technical ingredients now exist. Attention masking derived from panoptic masks provides a direct route from graph edges to spatial constraints on attention, and reuses the mechanism of Chapter~7 in a different role. The correspondence supervision of Chapter~6 generalizes naturally. Where \leffa{} supervises target-to-reference attention with a flow field, a graph-conditioned model can supervise entity-token-to-region attention with the panoptic mask of that entity, so that the assignment of tokens to entities is enforced rather than assumed. Synthetic supervision is also cheap here, because a graph can be extracted from any image with an existing PSG model, which yields arbitrarily many graph-image pairs without annotation. The price is that these pairs inherit the biases of that model.

\noindent \textbf{Evaluation challenges.}
Structure-based evaluation is attractive because it measures what current metrics cannot. One parses the generated image back into a scene graph and scores it against the graph that was requested, which measures constraint satisfaction directly rather than realism or aggregate semantic alignment. The obstacle is circularity. A parser that carries the biases of Chapter~3 is not a neutral judge, and a generator optimized against such a parser will learn its blind spots. Three mitigations deserve study: parser ensembles that use disagreement as an abstention signal, held-out parsers that are never used during training, and human adjudication on a stratified sample in which rare relations are deliberately over-represented. None is fully satisfactory, which makes the calibration work of Section~\ref{sec:future_open_world} a prerequisite for this program rather than a parallel concern.

\subsection{Open-World Structure with Calibration and Abstention}
\label{sec:future_open_world}
Section~\ref{sec:thesis_limitations} identifies the absence of abstention as a limitation shared by every system in this thesis, and Section~\ref{openpsg:sec:limitations} shows that it is most severe where a generative decoder produces fluent relations for pairs that may not be related. A system deployed in an open world encounters mostly unrelated entity pairs, so knowing when to say nothing is a precondition for open-vocabulary recognition rather than a refinement of it.

\noindent \textbf{Research questions.}
First, what is the right notion of confidence for a generated relation phrase? For a classifier, calibration is well defined over a fixed label set. For an autoregressive decoder, sequence likelihood conflates the plausibility of the relation with the fluency of its surface form, so a common paraphrase scores higher than a precise but unusual predicate. Whether a usable confidence signal can be recovered from generation remains open. Candidate sources include self-consistency across sampled decodings, agreement between generation and a separate verification pass, and an explicitly trained abstention head.

Second, can a model distinguish the case where no relation holds from the case where a relation holds that it cannot name? These are different states and require different responses, but no current PSG formulation represents the difference. Doing so would require a label space with an explicit null and an explicit unknown, together with training data in which both are annotated, and no existing dataset provides this.

Third, how should predicate vocabularies grow over time without catastrophic forgetting, and who decides what enters them? The second half of this question is not merely procedural. Section~\ref{sec:thesis_limitations} argues that making the supervised quantity explicit converts an implicit bias into an explicit one, and a vocabulary that expands through use makes that conversion continuous rather than one-off.

Fourth, how do relation distributions shift across domains, and can the shift be predicted? Long-tail predicates are domain-specific, so a robotics deployment and a consumer-photo deployment have different tails, and a model calibrated on one is miscalibrated on the other. The miscalibration may be systematic enough to correct.

\noindent \textbf{Methodological opportunities.}
Retrieval-augmented relation prediction suits the tail well. Instead of requiring the model to have internalized a rare predicate, the system retrieves visually or textually similar exemplars and conditions on them, so that rare-relation performance depends on the retrieval index rather than on the training distribution. This also turns vocabulary expansion into a data operation rather than a training operation. Conformal prediction offers distribution-free coverage guarantees and would produce relation \emph{sets} with calibrated confidence instead of point predictions, which suits a task where several predicates are simultaneously valid. That is precisely the observation from which Chapter~3 begins. The interaction-aware pair selection of \citet{li2025inova} also suggests that learning what makes a pair relation-bearing is more robust than thresholding a relation score, which is the filtering strategy that Chapter~5 uses.

\noindent \textbf{Evaluation challenges.}
The largest gap lies here, and Section~\ref{openpsg:sec:limitations} states the case. Current protocols measure held-out-vocabulary generalization while reporting it as open-set prediction, and exact-match scoring penalizes paraphrase. Progress requires benchmarks that admit semantic equivalence, whether through synonym sets, embedding-space matching, or reference-free image-relation alignment scoring~\citep{neau2025measuring}. It also requires benchmarks that reward abstention when evidence is genuinely absent, which means annotating pairs as unrelated rather than simply omitting them. Finally, it requires benchmarks that contain genuinely out-of-vocabulary relations with ground truth, which no current benchmark does. The last of these is the hardest and the most valuable. Building it would require annotators to write free-form relation phrases with no vocabulary to choose from, and then to reconcile the resulting variation. Reconciling that variation is itself a research problem, because the variation is the phenomenon such a benchmark exists to measure.

\subsection{Disentangling What a Reference Specifies}
\label{sec:future_reference}
Chapters~6 and~7 treat a reference as something to reproduce as faithfully as possible. Section~\ref{saber:sec:post_publication} identifies the assumption this conceals. A reference image mixes properties that are intrinsic to the subject with properties that are incidental to the moment of capture, and fidelity-maximizing training reproduces both. A prompt that asks for a change to an incidental property then conflicts directly with the training objective.

\noindent \textbf{Research questions.}
First, can intrinsic and incidental reference properties be separated without supervision? Identity, garment texture, and logo geometry should persist, whereas illumination, pose, camera characteristics, and background should not. No existing formulation makes this distinction, and it is unclear whether it can be learned from data alone or requires an explicit factorization.

Second, what should a model do when the reference and the prompt conflict? Current systems resolve such conflicts implicitly, according to whichever signal happens to dominate. An explicit priority mechanism, together with an interface through which a user states the priority, would make the behavior predictable. Predictability is arguably more valuable than any particular resolution.

Third, how should multiple references be fused when their relationship is unknown? Chapter~7 shows Saber degrading rather than saturating as references accumulate, because nothing distinguishes multiple views of one subject from views of several subjects. Inferring this grouping, rather than requiring it as input, is a well-posed and self-contained problem.

Fourth, does correspondence supervision transfer to diffusion transformers, and does it compose with bidirectional training? Section~\ref{leffa:sec:post_publication} notes that we validated \leffa{} on U-Net backbones and that Voost~\citep{lee2025voost} improves correspondence by a different route. Whether the two compound is an inexpensive experiment with a clear negative result available.

\noindent \textbf{Methodological opportunities.}
Cross-clip retrieval~\citep{xu2026vera} decorrelates the reference from the target photometrically, which the geometric augmentation of Saber does not do. The open question is whether the same decorrelation can be obtained without identity annotation, for example by pairing frames separated by long temporal gaps or by scene cuts, or by applying strong photometric augmentation such as relighting, color transfer, and synthetic sensor noise to the reference alone. Contrastive objectives over reference pairs offer a further route to disentanglement. Two references that depict the same subject under different conditions form a positive pair whose shared content is intrinsic by construction, which converts the disentanglement question into a data-construction question.

\noindent \textbf{Evaluation challenges.}
Current metrics reward fidelity and therefore cannot detect excessive fidelity. A model that copies incidental reference properties scores well on identity similarity while failing the prompt, and no standard metric registers this failure. Evaluation must therefore become conditional on intent. Given a prompt that requests a change to an incidental property, does the output make that change while preserving identity? Building such a benchmark requires prompts paired with references and an explicit annotation of which properties should vary. This is a modest annotation effort with a clear specification, and on the evidence of this thesis it would be unusually valuable.

\subsection{Long-Horizon Consistency and Memory}
The temporal structure studied in Chapter~7 spans a single clip. The harder regime is long-horizon generation, where a subject must remain consistent across shots, occlusions, and re-entries, and where errors accumulate rather than merely occur.

\noindent \textbf{Research questions.}
First, what is the right representation for persistent identity? Attention over reference tokens is adequate within a clip, but it does not obviously extend to a subject that leaves the frame and returns after a scene change. Whether an explicit memory, such as a bank of identity descriptors updated over time, outperforms a longer attention context, and at what horizon the crossover occurs, is unresolved.

Second, can identity drift be optimized against directly rather than only measured? A training objective that penalizes the divergence between a subject descriptor at frame $1$ and at frame $T$ would target the failure explicitly, in the same spirit as the correspondence supervision of Chapter~6. The obstacle is that such a descriptor is itself model-derived and can therefore be optimized in degenerate ways, since a model can minimize descriptor drift by generating a subject that barely moves.

Third, how should a system handle re-identification after occlusion or a cut? This is a tracking problem embedded in a generation problem, and it is unclear whether it is better solved by an explicit tracker or by an architecture in which the question does not arise.

Fourth, can generation be made interactive without becoming inconsistent? Users would rather correct one frame and have the correction propagate than regenerate an entire video. Propagation requires the correspondence machinery of Chapter~6 applied temporally, and its interactive form imposes latency constraints that current sampling budgets do not meet.

\noindent \textbf{Methodological opportunities.}
Hybrid architectures that combine diffusion with explicit state are the obvious candidate, whether through memory banks, state-space models, or recurrent latent summaries. The masked-training paradigm of Chapter~7 also extends naturally. Masking an entire temporal segment rather than a spatial region forces the model to carry identity across the gap, which makes long-horizon consistency a training objective rather than an emergent property.

\noindent \textbf{Evaluation challenges.}
Existing video metrics are computed per clip and aggregated over frames, so they cannot detect slow drift. A subject that changes gradually over thirty seconds may score well at every frame and at every adjacent pair. Metrics that compare distant frames directly, and benchmarks with clips long enough for drift to appear, are prerequisites. Multi-shot settings add a further difficulty, because consistency must be judged across cuts where viewpoint, lighting, and scale all change legitimately. Any such metric must therefore distinguish drift from intended variation, which returns to the disentanglement problem of Section~\ref{sec:future_reference}.

\subsection{Evaluation as a Research Object}
\label{sec:future_evaluation}
The recurring theme of Section~\ref{sec:thesis_limitations} is that this thesis measures agreement with annotations rather than correctness, and that the gap is unmeasured. This is not a limitation specific to this work. It is the state of both fields, and it deserves treatment as a research problem rather than as infrastructure.

\noindent \textbf{Research questions.}
First, how large is the gap between benchmark agreement and correctness, and can it be estimated? For PSG this is measurable in principle. One can exhaustively annotate a small set of images with every valid relation, then quantify how much standard recall understates performance and whether the understatement is uniform across the frequency spectrum. Chapter~3 predicts that it is not, and that omissions concentrate in the tail, so tail performance is understated most. Establishing this would recalibrate a decade of reported numbers.

Second, what is the inter-annotator agreement for relation annotation, and what does it imply for achievable performance? If annotators agree on only a fraction of the predicates for a given pair, then a model that matches one annotator exactly has reached a ceiling, and gains beyond that point measure conformity to a particular annotator rather than accuracy. \citet{li2024semantics} argue that annotator language preference is a principal source of the semantic overlap that Chapter~3 addresses, which suggests the ceiling is low enough to matter.

Third, can automatic metrics be validated against human judgment systematically rather than anecdotally? Chapter~6 reports FID alongside a human study and observes that the two disagree, and VTBench~\citep{hu2025vtbench} argues that this disagreement is structural. A study that measured the rank correlation between each standard metric and human preference, across many method pairs and stratified by failure type, would establish which metrics are safe to optimize. The answer may be that several widely used metrics are not.

Fourth, how should conflicting constraints be scored? When a prompt and a reference disagree, no single output is correct, so a metric that assumes a unique ground truth does not apply. Scoring the trade-off, that is, how gracefully a system degrades as constraints conflict, requires benchmarks built deliberately around conflict, and these do not exist.

\noindent \textbf{Methodological opportunities.}
Reference-free evaluation~\citep{neau2025measuring} removes the dependence on a fixed annotation vocabulary and is the most promising direction for open-vocabulary settings, although it substitutes a dependence on a judge model whose own biases must then be characterized. Decomposed benchmarks that score specific failure modes separately, such as the texture preservation, background consistency, and occlusion handling of VTBench, are more informative than aggregate scores and are also diagnostic, because they indicate which component to improve rather than only that something is wrong. Stratified human evaluation, in which annotators locate specific errors rather than state an overall preference, would have tested the claims of Chapter~6 far more sharply than the protocol we used, and the same applies to most human studies in this literature.

\noindent \textbf{Evaluation challenges.}
The recursion here is unavoidable. Evaluating evaluation requires ground truth about what good evaluation is, which is ultimately human judgment, which is itself variable and biased. The realistic goal is not a metric that is correct, but a set of metrics whose disagreements are understood, so that a claim supported by several metrics with uncorrelated failure modes carries more weight than a claim supported by a single aggregate score. This thesis would have been more convincing under such a regime, and the argument of Section~\ref{sec:unifying_argument} would have been testable rather than only illustrated.

\subsection{Efficiency, Adaptation, and Trustworthy Generation}
Two shorter directions follow from concerns raised above rather than forming programs in themselves.

\noindent \textbf{Efficiency and adaptation.}
Section~\ref{sec:thesis_limitations} observes that this thesis relocates cost rather than reducing it. Parameter-efficient tuning, modular conditioning branches that compose without joint retraining, and caching strategies all bear on deployment~\citep{zhang2023adding, wu2024towards, zhou2025openpsg}. The central research question is compositional: when independently trained structural objectives or conditioning modules are combined, do their effects add or interfere? Chapter~6 shows that one correspondence objective transfers across backbones, but nothing establishes that two such objectives, or a spatial objective combined with a temporal one, remain compatible. Persistent interference would limit the viability of modular adaptation.

\noindent \textbf{Trustworthy generation.}
As reference-based control improves, generating convincing imagery of real people becomes easier, and the methods of Chapters~6 and~7 apply directly to that purpose. Provenance tracking, watermarking, and edit auditing are the standard responses, and the research question is whether they can be integrated without degrading controllability or introducing brittle failures under benign edits. This thesis raises a more specific question. The mechanisms developed here make models better at preserving identity, which means they also make misuse easier in exactly the dimension that matters most. Whether identity preservation can be made conditional, so that it is faithful for consented references and degraded otherwise, is technically open and, given the direction of the field, increasingly urgent.

\section{Closing Remarks}
The central conclusion of this thesis is that reliable constraint satisfaction requires structure-aligned learning: task-relevant relationships and invariances must directly shape model architecture and representation, learning objectives and explicit supervision, decoding procedures, or scalable data construction and training strategies. In visual understanding, this structure comprises pixel-grounded entities and relations. In reference-conditioned generation, it comprises pose, spatial correspondence, and identity consistency over time. These structures are not identical, and the proposed methods do not reduce to one architecture. Their common role is to suppress shortcut solutions that reduce training loss without satisfying the intended task constraints.

This claim does not establish that any particular representation of structure is complete or unbiased. Several structures used here, including a fixed predicate vocabulary, a human pose skeleton, and a finite set of references, encode assumptions that may need to be replaced rather than refined. Nor does the thesis demonstrate an end-to-end system in which structure recovered by visual understanding directly constrains and evaluates generation; Section~\ref{sec:future_structure_conditioned} presents that step as future work. The five technical chapters establish a narrower conclusion: across two distinct visual settings, reliability improves when the relationships that define the task are made operational in learning.

We hope that these methods and perspectives provide a foundation for future work on structure-aligned learning across visual understanding and generation.
